\documentclass{article}
\usepackage{amsmath,amssymb,amsthm}
\newtheorem{proposition}{Proposition}
\begin{document}
\section*{A coarse-observation test for Q's action rewards}

Q, in ``Mechanisms and revelation,'' defines a reward \(r(\hat t,a,\theta)\)
on actual action and state. Its ``Dimensions of single-agent mechanisms''
observes that rewards can only respond to what an evaluator identifies, but
does not give a cycle characterization for a coarse action measurement. P,
in ``Conditional limit of the wrist frame,'' supplies a candidate physical
measurement: the wrist, motor encoders, and action history can leave the
relative object pose \(G\) unresolved. P explicitly does not establish that
distinct \(G\) values occur with identical full histories on the hardware.

The following finite, deterministic measurement model isolates the incentive
issue. Fix one Q type, with finite signal set \(S\), feasible physical action
set \(A\), and posterior intrinsic utility \(U(a,s)\). Interpret an action as
an achieved physical trajectory when that interpretation is justified.
Suppose the contractible measurement is \(m:A\to Y\), and the reward must
factor as \(r(a)=q(m(a))\), where \(q\) may prohibit an observation label by
assigning it \(-\infty\). For a target policy \(a_s\in A\), put
\[
Y_0=\{m(a_s):s\in S\}.
\]
Construct a directed multigraph on \(Y_0\). For every \(s\in S\) and
\(b\in A\) with \(m(b)\in Y_0\), insert an edge from \(m(a_s)\) to
\(m(b)\) of weight \(U(b,s)-U(a_s,s)\).

\begin{proposition}
The target policy has a supporting observation-based reward if and only if
every directed cycle in this graph has nonpositive total weight.
\end{proposition}

\begin{proof}
Obedience requires
\[
q(m(a_s))-q(m(b))\geq U(b,s)-U(a_s,s).
\]
Summing these inequalities around a graph cycle proves necessity. If there
is no positive cycle, let \(q(y)\) be the maximum weight of a path starting
at \(y\), including the empty path. Finite nodes and nonpositive cycles
make this finite. Prepending an edge of weight \(w\) gives
\(q(y)\geq w+q(y')\), the obedience inequality. Set \(q(y)=-\infty\)
outside \(Y_0\). Every on-policy label remains available, and all
off-policy labels are prohibited.
\end{proof}

A positive self-loop is already decisive: if \(m(b)=m(a_s)\) but
\(U(b,s)>U(a_s,s)\), no payment magnitude can make \(a_s\) optimal.
This extends Q's signal-policy cycle test by identifying actions that the
evaluator cannot distinguish. Q's more elaborate type-report constraints
would still need to be checked after each type's observation-based
supporting schedule is found. The statement assumes that the physical
alternatives are genuinely feasible for the agent and that their
measurements are deterministic. It does not follow merely from P's visible
in-hand rotation, since P reports no matched-history pose measurements and
models no strategic donor.

Separately, P's receiver-side information limit has a quantitative form.
For two equally likely object orientations \(H\in\{0,1\}\), let the receiver
observe \(Y\) with conditional laws \(P_0,P_1\). If no receiving action
succeeds under both orientations, any policy based only on \(Y\) has average
success at most
\[
\frac{1+\operatorname{TV}(P_0,P_1)}{2}.
\]
The proof is the binary classification bound: the policy must choose one of
two mutually exclusive success sets, and its optimal classification
probability is the displayed quantity. Any Q-style messages or rewards
generated solely from \(Y\) and independent randomization cannot increase
the bound, by data processing. A donor's genuinely pose-sensitive private
signal could increase it; the relevant observation would then include that
signal or a truthful report of it. A two-point experiment consisting only of
P's hypothesized matched trials has \(P_0=P_1\), giving \(1/2\). A single
matched pair does not show equality of observation laws in routine trials,
and P does not establish such equality empirically.

The pair suggests a testable question: measure how the object pose within
the wrist varies across trials with matched contractible histories, then
check whether the resulting observation graph contains a profitable hidden
deviation for a specified donor objective, and whether adding a pose-sensitive
sensor changes receiver success. The graph condition is a generic finite
difference-constraint fact, so it is not by itself a novelty claim.
\end{document}
